\documentclass[11pt,a4paper]{article}
\usepackage[margin=2.2cm]{geometry}
\usepackage{amsmath,amssymb,booktabs,graphicx,enumitem,xcolor,array,float}
\usepackage[hidelinks]{hyperref}
\usepackage[font=small,labelfont=bf]{caption}
\setlength{\parskip}{4pt}\setlength{\parindent}{0pt}
\newcommand{\TODO}[1]{\textcolor{red}{[#1]}}
\ifdefined\PublicBuild% Team details, shared by the proposal, the progress report and the final report.
% Fill in the names once here. For the printed submission, copy this file to team_private.tex
% (git-ignored) and add student IDs there; the reports use team_private.tex when it exists.
\providecommand{\TODO}[1]{\textcolor{red}{[#1]}}
\newcommand{\TeamNames}{Nguyen Van Thanh, \TODO{Member 2}, \TODO{Member 3}, \TODO{Member 4}, \TODO{Member 5}}
\newcommand{\TeamNamesStacked}{Nguyen Van Thanh\\ \TODO{Member 2} \quad \TODO{Member 3} \quad \TODO{Member 4} \quad \TODO{Member 5}}
% owner of each method (one method per member)
\newcommand{\MOne}{\TODO{Member A}}
\newcommand{\MTwo}{\TODO{Member B}}
\newcommand{\MThree}{\TODO{Member C}}
\newcommand{\MFour}{\TODO{Member D}}
\newcommand{\MFive}{\TODO{Member E}}
\newcommand{\SubmissionDate}{\TODO{Submission date}}
\else\IfFileExists{../team_private.tex}{\input{../team_private.tex}}{% Team details, shared by the proposal, the progress report and the final report.
% Fill in the names once here. For the printed submission, copy this file to team_private.tex
% (git-ignored) and add student IDs there; the reports use team_private.tex when it exists.
\providecommand{\TODO}[1]{\textcolor{red}{[#1]}}
\newcommand{\TeamNames}{Nguyen Van Thanh, \TODO{Member 2}, \TODO{Member 3}, \TODO{Member 4}, \TODO{Member 5}}
\newcommand{\TeamNamesStacked}{Nguyen Van Thanh\\ \TODO{Member 2} \quad \TODO{Member 3} \quad \TODO{Member 4} \quad \TODO{Member 5}}
% owner of each method (one method per member)
\newcommand{\MOne}{\TODO{Member A}}
\newcommand{\MTwo}{\TODO{Member B}}
\newcommand{\MThree}{\TODO{Member C}}
\newcommand{\MFour}{\TODO{Member D}}
\newcommand{\MFive}{\TODO{Member E}}
\newcommand{\SubmissionDate}{\TODO{Submission date}}
}\fi
\newcommand{\fig}[1]{../../figures/#1}

\begin{document}
\begin{center}
{\large Operations Research 76000093 --- Final Project Progress Report}\\[6pt]
{\LARGE\bfseries Right Place, Right Time: Capacity, Guidance and Pricing\\[2pt] for Public EV Charging in Shenzhen}\\[8pt]
\TeamNames\qquad 30 November 2026
\end{center}

\section{The problem we are studying, and why}
Shenzhen had 1.08 million new-energy vehicles by June 2024~\cite{cnev2024}, and the national NEV stock grew by 49\% in the following year~\cite{cnev2025}. We study whether the city's \textbf{17{,}532 public charging piles in 275 traffic zones} (UrbanEV dataset~\cite{li2025urbanev}, hourly, Sep 2022--Feb 2023) can absorb this growth, and how to fix it at least cost. The data show a paradox. Average occupancy is only 28\%, yet a few zones saturate, mostly at night. Our working hypothesis, confirmed so far, is that the network suffers from a \emph{spatial and temporal mismatch} rather than a simple shortage. This matters because the cheap levers (moving idle piles, guiding drivers, re-timing prices) could replace much of the expensive one (building new piles).

\textbf{Decision problem.} Keep driver waiting time low under demand growth $g \in \{1.5, 2, 2.5\}$, choosing among four levers: new piles, relocation of existing piles, navigation guidance, and a time-of-use service fee.

\section{Methods and why we chose them}
\begin{center}\small
\begin{tabular}{@{}p{2.4cm}p{5.6cm}p{7.6cm}@{}}\toprule
Method & What it answers & Why this method\\\midrule
M1 Queueing (M/M/$c$) & waiting time per zone and hour; which zones fail as demand grows & Charging piles are parallel servers and waiting is the service measure drivers feel; Erlang-C links piles, load and delay in closed form and needs only the data we have (busy piles).\\
M2 Integer programming & where to add and where to remove piles under a budget & Piles are indivisible and the budget couples all zones. The waiting cost is convex in piles, so an exact 0-1 incremental MILP is possible, and LP duality gives a shadow price of the budget.\\
M3 Network flow & how much overload guidance can absorb within a detour limit & Redirecting drivers is a transportation problem (supplies = excess, demands = spare capacity, costs = distance); total unimodularity guarantees integer flows and fast solution (network simplex).\\
M4 Nonlinear programming & which hourly fee shifts demand away from congested hours & Demand response to price is nonlinear (logit), and revenue neutrality is a nonlinear equality; KKT multipliers quantify the trade-off between revenue and congestion.\\
M5 Simulation + simulated annealing & do M1/M2 hold under time-varying demand and spill-over between zones? & The analytical models assume stationarity and independent zones; only simulation captures interactions. The resulting objective is noisy and non-separable, which suits a metaheuristic.\\\bottomrule
\end{tabular}\end{center}

\section{How we accomplish the work (pipeline)}
\begin{enumerate}[leftmargin=1.5em,itemsep=1pt]
\item \textbf{Data}: per-zone pile counts, hourly offered load (mean day), zone distances, prices. We found that \texttt{occupancy.csv} holds counts, not percentages as documented, and verified this against pile counts.
\item \textbf{Demand}: the mean day is not censored, while busy days are (10 saturated zones), so we use the mean day $\times\,g$ and study variability by simulation.
\item \textbf{Models}: M1 $\rightarrow$ waiting curves $W_z(c)$ feed M2. M3 and M4 produce modified loads, which are fed back into M1/M2. M5 re-evaluates M2-type plans.
\item \textbf{Verification}: each model is checked independently: Erlang formulas against textbook values, MILP against its LP relaxation, LP against network simplex, NLP solvers against each other with a KKT residual check, simulation against Erlang-C.
\item \textbf{Common yardstick}: the minimum budget needed to reach an average wait below 15\,s per driver, to compare all levers.
\end{enumerate}

\section{Preliminary results (all for doubled demand, $g = 2$)}
\begin{itemize}[leftmargin=1.3em,itemsep=2pt]
\item \textbf{M1}: 48 zones become unstable although network occupancy is only 56\%; the piles needed for stability number 598. A useful theoretical result: total waiting time does not depend on the unknown mean session length ($\lambda W_q = aC/(c-a)$).
\item \textbf{M2}: relocation dominates. At a budget of 1{,}000 new-pile units, the optimal mix of new and relocated piles leaves 5.0 driver-hours of waiting per day, against 952 with new piles only; a proportional rule leaves 41 zones unstable. The shadow price of the budget falls from 105 to 0.01 driver-hours/day per unit between $B = 250$ and $1000$.
\item \textbf{M3}: guidance within 3\,km redirects 70\% of the excess load and cuts the piles needed for stability from 598 to 185. LP and network simplex agree (748.1 vs 745.3 at the peak hour, within discretisation).
\item \textbf{M4}: the estimated price sensitivity is $\hat\beta = 0.281$ per yuan/kWh ($t = -2.21$ for the slope $-\beta$). The revenue-neutral optimal fee halves the congestion index but changes the number of unstable zones only from 48 to 45: price mostly helps in time, not in space.
\item \textbf{M5}: the simulation reproduces Erlang-C (all test cases within 0.35 standard errors). With driver spill-over, simulated annealing improves on the Erlang-optimal plan by 34\% out of sample.
\end{itemize}
\begin{figure}[H]\centering
\includegraphics[width=0.49\linewidth]{\fig{m2_budget_curve.png}}\hfill
\includegraphics[width=0.49\linewidth]{\fig{m3_guidance_curve.png}}
\caption{Preliminary results: M2 waiting vs budget (left) and M3 effect of the detour limit (right).}\end{figure}

\section{Remaining work, schedule and task allocation}
\begin{center}\small
\begin{tabular}{@{}p{3cm}p{12.6cm}@{}}\toprule
Period & Tasks\\\midrule
Dec 1--7 & Integrated policy analysis (all levers on the budget yardstick); robustness of pricing (moderate fee range); growth sensitivity for M2--M4.\\
Dec 8--14 & Literature review write-up (10 papers: summary, critique, gaps); difficulties section; recommendations.\\
Dec 15--(TBA) & Final report integration; 10--15 slides with one section per member; two full rehearsals with timing (10--12 min).\\\bottomrule
\end{tabular}\end{center}
\begin{center}\small
\begin{tabular}{@{}p{3.4cm}p{12.2cm}@{}}\toprule
Member & Responsibility (status)\\\midrule
\MOne & M1 queueing, data pipeline, growth scenarios (done; writing up)\\
\MTwo & M2 MILP, duality and sensitivity (done; integrated analysis in progress)\\
\MThree & M3 min-cost flow, network simplex (done; combination with M4 in progress)\\
\MFour & M4 estimation and NLP (done; robustness of fee ranges in progress)\\
\MFive & M5 simulation and simulated annealing (done; sensitivity to spill-over parameters planned)\\
All & literature review (two papers each), final report, presentation\\\bottomrule
\end{tabular}\end{center}
\textbf{Risks and mitigation.} (i) Weak identification of $\beta$: we report results for $\beta \in [0.1, 1]$. (ii) Censored demand underestimates today's shortage: we plan with growth scenarios and state this as a limitation. (iii) Five members (approved by the instructor) instead of 3--4: we cover five distinct methods, one per member.

\begin{thebibliography}{9}\small
\bibitem{li2025urbanev} H.~Li, H.~Qu, X.~Tan, L.~You, R.~Zhu, W.~Fan, ``UrbanEV: An open benchmark dataset for urban electric vehicle charging demand prediction,'' \emph{Scientific Data}, 12:523, 2025. doi:10.1038/s41597-025-04874-4.
\bibitem{cnev2024} CnEVPost, ``Shenzhen NEV penetration reaches record 77.4\% in Jun,'' Jul.~18, 2024.
\bibitem{cnev2025} CnEVPost, ``China NEV ownership reaches 36.89 million,'' Jul.~14, 2025.
\end{thebibliography}
\end{document}
